% =============================================================================
% CONCLUSION
% =============================================================================

\label{sec:conclusion}

This paper introduces CITA (Contrastive Instruction-Tuned Alignment), a novel framework that enables instruction-aware alignment in large language models. Unlike traditional alignment methods (RLHF, DPO) that embed static alignment objectives into model weights, CITA allows models to dynamically adapt their behavior based on alignment instructions provided at inference time.

\subsection{Key Findings}

\begin{enumerate}
    \item \textbf{CITA outperforms DPO in instruction-awareness}: When comparing the improvement from NoInstruct to Instruct variants, CITA wins 3 out of 5 evaluation metrics (TruthfulQA, Length Control, AQI), while DPO wins 2 (ISD, Conditional Safety).

    \item \textbf{Correct directional adaptation}: On TruthfulQA, CITA improves by +0.054 (becoming more confident when instructed), while DPO shows minimal change (+0.001), indicating CITA better follows uncertainty instructions.

    \item \textbf{Mandatory KL regularization is essential}: The KL term prevents catastrophic forgetting and enables stable instruction-conditioned behavior switching.

    \item \textbf{Instruction-alignment $\neq$ instruction-following}: Through the Instruction Switch Dataset, we demonstrate that CITA achieves deeper instruction-alignment rather than superficial instruction-following.
\end{enumerate}

\subsection{Implications}

CITA enables a new paradigm of \textbf{interactive and adaptive AI systems} where:
\begin{itemize}
    \item Policy updates can be expressed in natural language without retraining
    \item Different users or contexts can receive appropriately calibrated responses
    \item Safety policies can be dynamically adjusted based on deployment requirements
\end{itemize}

